\documentclass{article}
\usepackage{graphicx} % Required for inserting images

\title{s2_g3_ITS_L9_Scribe_CoT_AU2440036_Virja}
\author{Virja Shah}
\date{February 2026}

\usepackage{amsmath, amssymb}

\title{Lecture 9: Uniform, Exponential, Laplace and Gamma Random Variables}
\author{}
\date{}

\begin{document}
\maketitle

\section*{1. Outline of Lecture (from slides)}
This lecture covers Types of Continuous Random Variables, specifically:
\begin{itemize}
    \item Uniform Random Variable
    \item Exponential Random Variable
\end{itemize}

\noindent (Laplace and Gamma are mentioned in the title/outline but this particular provided PDF segment contains Uniform + Exponential content only.)

% ----------------------------

\section*{2. Uniform Random Variable}

\subsection*{2.1 Definition (PDF)}
A continuous random variable $X$ is said to be Uniform over the interval $[a,b)$ if its PDF is:

\[
f_X(x)=\frac{1}{b-a} \quad \text{for } a\le x<b
\]
\[
f_X(x)=0 \quad \text{otherwise}
\]

\noindent Interpretation:
\begin{itemize}
    \item All values between $a$ and $b$ are equally likely.
    \item The height of the PDF is constant: $\frac{1}{b-a}$.
\end{itemize}

\subsection*{2.2 Definition (CDF)}
The CDF of a uniform random variable is:

\[
F_X(x)=0 \quad \text{for } x<a
\]
\[
F_X(x)=\frac{x-a}{b-a} \quad \text{for } a\le x<b
\]
\[
F_X(x)=1 \quad \text{for } x\ge b
\]

\noindent Interpretation:
\begin{itemize}
    \item For values before $a$, probability accumulated is 0.
    \item Between $a$ and $b$, probability increases linearly.
    \item After $b$, probability becomes 1.
\end{itemize}

\subsection*{2.3 Graph (from slide figure)}
From the diagram:

\noindent (a) PDF
\begin{itemize}
    \item A flat horizontal line between $a$ and $b$ at height $\frac{1}{b-a}$.
    \item 0 everywhere else.
\end{itemize}

\noindent (b) CDF
\begin{itemize}
    \item 0 until $x=a$.
    \item Then a straight line increasing from 0 to 1.
    \item Then constant at 1 after $x=b$.
\end{itemize}

% ----------------------------

\section*{3. Example 1 (Uniform Random Variable)}

\subsection*{3.1 Problem Statement (given)}
The phase of a sinusoid, $\Theta$, is uniformly distributed over $[0,2\pi)$.  
So its PDF is:

\[
f_\Theta(\theta)=\frac{1}{2\pi} \quad \text{for } 0\le \theta<2\pi
\]
\[
f_\Theta(\theta)=0 \quad \text{otherwise}
\]

\subsection*{Key Uniform Property Used (from solution slide)}
For a uniform RV on $[0,2\pi)$:

\[
\Pr(a<\Theta<b)=\frac{b-a}{2\pi}
\]

\subsection*{(a) Find $\Pr\left(\Theta>\frac{3\pi}{4}\right)$}
Total interval length $=2\pi$.

\noindent Favorable region $=\left(\frac{3\pi}{4},2\pi\right)$.

\[
\Pr\left(\Theta>\frac{3\pi}{4}\right)=\frac{2\pi-\frac{3\pi}{4}}{2\pi}
\]

Compute numerator:

\[
2\pi-\frac{3\pi}{4}=\frac{8\pi}{4}-\frac{3\pi}{4}=\frac{5\pi}{4}
\]

Thus:

\[
\Pr\left(\Theta>\frac{3\pi}{4}\right)=\frac{\frac{5\pi}{4}}{2\pi}=\frac{5}{8}
\]

\subsection*{(b) Find $\Pr\left(\Theta<\pi\mid \Theta>\frac{3\pi}{4}\right)$}
Let:

\[
A=\{\Theta<\pi\}
\]
\[
B=\left\{\Theta>\frac{3\pi}{4}\right\}
\]

We use conditional probability:

\[
\Pr(A\mid B)=\frac{\Pr(A\cap B)}{\Pr(B)}
\]

From part (a):

\[
\Pr(B)=\frac{5}{8}
\]

Now compute intersection:

\[
A\cap B=\left\{\frac{3\pi}{4}<\Theta<\pi\right\}
\]

Using uniform interval property:

\[
\Pr\left(\frac{3\pi}{4}<\Theta<\pi\right)=\frac{\pi-\frac{3\pi}{4}}{2\pi}
\]

Compute numerator:

\[
\pi-\frac{3\pi}{4}=\frac{4\pi}{4}-\frac{3\pi}{4}=\frac{\pi}{4}
\]

So:

\[
\Pr(A\cap B)=\frac{\frac{\pi}{4}}{2\pi}=\frac{1}{8}
\]

Now substitute:

\[
\Pr(A\mid B)=\frac{\frac{1}{8}}{\frac{5}{8}}=\frac{1}{5}
\]

\subsection*{(c) Find $\Pr\left(\cos(\Theta)<\frac{1}{2}\right)$}
We use the trig identity given in slide:

\[
\cos\Theta=\frac{1}{2}\Rightarrow \Theta=\frac{\pi}{3},\frac{5\pi}{3}
\]

Cosine is less than $\frac{1}{2}$ in the interval:

\[
\frac{\pi}{3}<\Theta<\frac{5\pi}{3}
\]

So:

\[
\Pr\left(\cos(\Theta)<\frac{1}{2}\right)=\frac{\frac{5\pi}{3}-\frac{\pi}{3}}{2\pi}
\]

Compute numerator:

\[
\frac{5\pi}{3}-\frac{\pi}{3}=\frac{4\pi}{3}
\]

Thus:

\[
\Pr\left(\cos(\Theta)<\frac{1}{2}\right)=\frac{\frac{4\pi}{3}}{2\pi}=\frac{4}{6}=\frac{2}{3}
\]

% ----------------------------

\section*{4. Uniform Random Variable: Application Examples}
The slides list these direct applications:

\begin{itemize}
    \item Phase of a sinusoidal signal, where all phase angles between $0$ and $2\pi$ are equally likely.
    \item Random number generated by a computer between 0 and 1 for simulations.
    \item Arrival time of a user within a known time window, assuming no time preference.
\end{itemize}

% ----------------------------

\section*{5. Exponential Random Variable}

\subsection*{5.1 PDF Definition (given)}
The exponential random variable has PDF (for any $b>0$):

\[
f_X(x)=\frac{1}{b}\exp\left(-\frac{x}{b}\right)u(x)
\]

\noindent Here $u(x)$ is the unit step function.  
This implies the PDF is only valid for $x\ge 0$.

\subsection*{5.2 CDF Definition (given)}
The exponential random variable has CDF:

\[
F_X(x)=\left[1-\exp\left(-\frac{x}{b}\right)\right]u(x)
\]

Again, $u(x)$ ensures the distribution is only for $x\ge 0$.

\subsection*{5.3 Graph Interpretation (from slides)}
From the plotted graphs:

\noindent PDF:
\begin{itemize}
    \item Starts at $\frac{1}{b}$ when $x=0$
    \item Decreases exponentially toward 0 as $x\to\infty$
\end{itemize}

\noindent CDF:
\begin{itemize}
    \item Starts at 0 when $x=0$
    \item Increases toward 1 as $x\to\infty$
\end{itemize}

% ----------------------------

\section*{6. Example 2 (Exponential Random Variable)}

\subsection*{6.1 Problem Statement}
Let $X$ be an exponential random variable with PDF:

\[
f_X(x)=e^{-x}u(x)
\]

\subsection*{Identify parameter $b$}
General exponential PDF from slide:

\[
f_X(x)=\frac{1}{b}\exp\left(-\frac{x}{b}\right)u(x)
\]

Given:

\[
f_X(x)=e^{-x}u(x)
\]

So we match:

\[
\frac{1}{b}=1 \Rightarrow b=1
\]
\[
-\frac{x}{b}=-x \Rightarrow b=1
\]

Thus:

\[
b=1
\]

\subsection*{(a) Find $\Pr(3X<5)$}
\[
3X<5 \Rightarrow X<\frac{5}{3}
\]

Now use the CDF formula from slides:

For $b=1$:

\[
F_X(x)=1-e^{-x}
\]

Thus:

\[
\Pr\left(X<\frac{5}{3}\right)=F_X\left(\frac{5}{3}\right)
\]

So:

\[
\Pr(3X<5)=1-e^{-5/3}
\]

\subsection*{(b) Generalize to find $\Pr(3X<y)$}
\[
3X<y \Rightarrow X<\frac{y}{3}
\]

So:

\[
\Pr(3X<y)=\Pr\left(X<\frac{y}{3}\right)=F_X\left(\frac{y}{3}\right)
\]

Using:

\[
F_X(x)=1-e^{-x}
\]

Therefore:

\[
\Pr(3X<y)=1-e^{-y/3}
\]

% ----------------------------

\section*{7. Final Exam Revision Summary (Lecture 9)}

\subsection*{Uniform RV}

PDF:
\[
f_X(x)=\frac{1}{b-a}, \quad a\le x<b
\]

CDF:
\[
F_X(x)=\frac{x-a}{b-a}, \quad a\le x<b
\]

Probability in interval:
\[
\Pr(a<X<b)=\frac{b-a}{\text{total length}}
\]

\subsection*{Exponential RV}

PDF:
\[
f_X(x)=\frac{1}{b}e^{-x/b}u(x)
\]

CDF:
\[
F_X(x)=\left[1-e^{-x/b}\right]u(x)
\]

\end{document}

